\begin{abstract}
Weight-space model merging has emerged as a computationally efficient paradigm for consolidating task-specific expert neural networks without the prohibitive costs of multi-task retraining. 
However, when ensembling coefficients are dynamically optimized on extremely scarce calibration data (e.g., ten samples per task), the optimizer is vulnerable to severe transductive overfitting, which we term the Overfitting-Optimizer Paradox. 
To resolve this, we present \textbf{PAC-Bayes Merge}, the first formal statistical learning-theoretic framework for trajectory-regularized model merging. 
By projecting layer-wise ensembling coefficients onto a low-degree polynomial trajectory and modeling them as the mean parameters of a randomized Gaussian posterior classifier, we prove that minimizing Alquier's linear PAC-Bayesian generalization bound analytically yields a quadratic $L_2$ Consensus-Pulling penalty centered at the stable, uniform ensembling consensus. 
Unlike heuristic $L_1$ penalties that force artificial coordinate sparsity and flatten trajectories, our smooth $L_2$ regularization preserves the continuous representative capacity of intermediate layers while robustly suppressing high-frequency validation noise. 
Evaluating on genuine physical image classification datasets (MNIST, FashionMNIST, CIFAR-10, SVHN) projected via Johnson-Lindenstrauss embeddings into a 14-layer deep MLP residual backbone across 15 random seeds, our advanced Fisher-guided non-isotropic PAC-Bayes-FIM Merge successfully circumvents transductive overfitting, yielding a Joint Mean accuracy of \textbf{36.13\%} under extreme data scarcity, outperforming the Static Uniform baseline (33.57\%), properly tuned Ties-Merge (29.68\%), DARE-Merge (33.24\%), and unconstrained layer-wise tuning (36.09\%). 
Furthermore, our randomized training and posterior ensemble evaluation bridge the theory-to-practice gap, providing rigorous learning-theoretic guarantees and high-performance deep learning.
\end{abstract}
